\pagenumbering{Roman}
\setcounter{page}{1}

\begin{center}
    {\heiti\zihao{-3} 拜占庭容错的多智能体优化} \\[2ex]   % 论文题目小3黑体
    {\heiti\zihao{3} 摘要} \\[2ex]                         % “摘要”3号黑体
    {\normalfont\zihao{-4} 学生：王梓恒 \qquad 指导教师：刘雪峰} \\[2ex]  % 请替换实际姓名
\end{center}

\noindent \quad \quad  分布式多智能体优化旨在通过无中心节点间的局部协作，最小化所有节点局部代价函数的算术平均值目标。然而，传统的算术平均优化目标在面临拜占庭故障时很脆弱，使得哪怕单一的恶意节点也能通过注入极端或不一致数据完全操控全局的优化结果。为了解决该问题，本文重构全局目标为仅由正常节点局部函数构成的凸组合。为了防止正常节点的数据权重被恶意节点稀释，本文首次提出了($\beta, \gamma$)-可容许性这一数学度量，用于限制正常节点占比不小于$\beta$的个数不小于$\gamma$。此外，本文通过聚焦标量局部代价函数，迈出了解决所提出的拜占庭容错多智能体优化问题的第一步。本文的研究结果可为一般化的局部代价函数研究提供有用的见解。

\vspace{1\baselineskip}
\noindent {\heiti\zihao{4} 关键词}：{\zihao{-4} 计算与人工智能，容错性，容错控制，机器学习/分布式计算，可靠性}
